\scanpage{2}
% archive foliation: 1
Connexions projectives canoniques et uniformisation des courbes algébriques.

1. Connexions projectives

 a) $S$ schéma de base (disons). La donnée d'un fibré de Brauer-Severi \inserted{$P\deleted{\unsure{S}}$} sur $S$ (schém. propre et lisse/$S$ dont les fibres \inserted{géom.} sont des espaces projectifs) de dim relative \inserted{constante} $r$ revient à la donnée d'un torseur $\mathcal{T}$ sur $S$ de Groupe $\mathrm{GP}(r)_S$. La donnée d'un $P$ avec section marquée revient à se donner une restriction du groupe structural au sous-Groupe \inserted{de $\mathrm{GP}(r)_S$} qui fixe une section donnée de \deleted{\unsure{$\mathcal{O}$}} $P_0^r=\mathbb{P}(\mathcal{O}_S^{r+1})$, disons celle qui correspond au quotient $\mathcal{O}_S$ de $\unsure{\mathcal{O}_S^{r+1}}$ ``correspondant à la dernière coordonnée''. Ce sous-groupe s'identifie au groupe affine $\mathrm{Aff}(r\deleted{\unsure{+1}})_S$, produit \unsure{tordu}
\[
\mathrm{Aff}(r\deleted{\unsure{+1}})=\mathrm{GL}(r\deleted{\unsure{+1}}).\mathbb{E}^r
\]
où $\mathrm{GL}(r)$ opère sur $\mathbb{E}^r$ de la façon évidente.

[\unsure{les|Le} \inserted{Groupe des} automorphismes du Module $\mathcal{O}_S^{r+1}=\mathcal{O}_S^r\times\mathcal{O}_S$ \inserted{$=E_0\times L_0$} qui respectent la structure d'extension, modulo les $\unsure{\mathrm{mult}^{\mathrm{ons}}}$ scalaires, \deleted{\unsure{est}} est en fait celui des matrices
\diagram{\[
{}_{r+1}\left\{\left(\begin{array}{c|c}
\underbrace{A}_{r}&C\\ \hline
0&\unsure{1}
\end{array}\right)\right.\!\begin{array}{l}\big\}r\\{}\end{array}
\]}
avec $A$ \deleted{\unsure{scalaires}} dans $\mathrm{GL}(r)$, $C$ dans $\mathbb{E}^r$ ---]

On trouve aussi que la donnée de $(P,\sigma)$ \uline{équivaut} à celle d'une \inserted{structure d'}extension
\deleted{de}
